On considère la fonction $f$ définie sur $\R$ par 

\[f(x)=5 \ln \left(x^{2}+1\right)-3 x\]
et on admet que la fonction $f$ est dérivable sur $\R$.

On note $\mathcal{C}_{f}$ sa courbe représentative dans un repère orthonormé du plan.

On a tracé ci-dessous la courbe $\mathcal{C}_{f}$ et la tangente à la courbe $\mathcal{C}_{f}$ au point A d'abscisse 1.

\begin{center}
\begin{tikzpicture}%compiler avec pgfplots
\begin{axis}[
	x={9mm}, y={9mm},
	xmin = -6.5, xmax = 6.9,
	ymin = -2.7, ymax = 6.1,
	xtick distance=1,
	ytick distance=1,
	minor tick num = 0,
	grid=none, axis lines =center]
	\node[below left=3pt] at (axis cs: 0,0) {{\footnotesize 0}};
	\addplot [line width=1.2pt, color=red, smooth, samples=100, domain= -1.5:6.9 ]{5*ln(x*x+1)-3*x};
	\node[red] at(axis cs: -1,3) {$\mathcal{C}_f$};
	\addplot [line width=1pt, color=teal, smooth, samples=10, domain= -1.5:6.9 ]{2*x-1.534};
	\draw[shift={(axis cs: 1,0.466)}] (-2pt,-2pt)--(2pt,2pt) (-2pt,2pt)node[above left]{A}--(2pt,-2pt) ;
\end{axis}
\end{tikzpicture}
\end{center}
%Figure en PSTricks
%\begin{center}
%\psset{unit=1cm,arrowsize=2pt 3}
%\begin{pspicture*}(-6,-2.5)(7,6)
%\psaxes[linewidth=1.25pt,labelFontSize=\scriptstyle](0,0)(-5.9,-2.5)(6.9,5.9)
%\psplot[plotpoints=2000,linewidth=1.25pt,linecolor=blue]{-2}{7}{x dup mul 1 add ln 5 mul 3 x mul sub}
%\psplotTangent{1}{4}{x dup mul 1 add ln 5 mul 3 x mul sub}
%\uput[ul](1,0.466){A}\psdots(1,0.466)
%\uput[l](-1,5){\blue $\mathcal{C}_f$}
%\end{pspicture*}
%\end{center}
%Fin de la figure en PSTricks

\begin{enumerate}
\item Conjecturer, à l'aide de la représentation graphique de la fonction $f$, les intervalles de $\R$ sur lesquels la fonction $f$ semble convexe ou concave.
\item Déterminer, en justifiant, la limite de la fonction $f$ en $-\infty$.
\item
	\begin{enumerate}
		\item Démontrer que, pour tout $x$ réel strictement positif,

\[f(x)=x\left(10 \dfrac{\ln x}{x}-3\right)+5 \ln \left(1+\dfrac{1}{x^{2}}\right)\]

		\item Déterminer, en justifiant, la limite de la fonction $f$ en $+\infty$.\end{enumerate}
\item
	\begin{enumerate}
		\item Démontrer que pour tout $x$ réel, \quad $f'(x)= \dfrac{-3 x^{2}+10 x-3}{x^{2}+1}$.
		\item Étudier les variations de la fonction $f$ sur $\R$.
	\end{enumerate}
	\item On admet que la fonction $f$ est deux fois dérivable sur $\R$ et que pour tout réel $x$,
	\[f''(x)=\dfrac{-10 x^{2}+10}{\left(x^{2}+1\right)^{2}}\]

	\begin{enumerate}
		\item Valider ou rejeter la conjecture faite à la question \textbf{1}.
		\item Déterminer l'équation réduite de la tangente à la courbe $\mathcal{C}_{f}$ au point A d'abscisse 1.
		\item En déduire que pour tout $x \geqslant 1$,\quad $\ln \left(x^{2}+1\right) \leqslant x+\ln (2)-1$.
	\end{enumerate}
	\end{enumerate}
